\newpage
\section{Approximate APSP with Additive Factors}\label{sc:additive}
In this section we give our second result, which is maintaining APSP on unweighted graphs with a mixed additive and multiplicative stretch. We start by reviewing useful tools in \Cref{sec:additive_tools}.  We then provide a decremental $(1+\epsilon, 2)$-APSP algorithm as a warm-up in \Cref{sec:additive_warm_up}. Finally in \Cref{sec:add_general} we generalize our techniques to maintain a $(1+\epsilon, 2(k-1))$-APSP algorithm efficiently.


\subsection{Decremental Tools}  \label{sec:additive_tools}

We start with describing several useful tools we use in our algorithms.

\remove{
\paragraph{Even-Shiloach tree (ES-tree).} An ES-tree is a central data structure for maintaining distances in a decremental setting. It maintains the exact distances from a source $s$ under edge deletions, giving the following.

\begin{lemma}(\cite{ES81,henzinger1995fully,king1999fully})\label{lemma_ES}
There is a data structure called ES-tree that, given a weighted directed graph $G$ undergoing deletions and edge weight increases, a root node $s$, and a depth parameter $d$, maintains, for every node $v$ a value $\ell_s(v)$ such that $\ell_s(v) = d_G(s,v)$ if $d_G(s,v) \leq d$ and $\ell_s(v) = \infty$ if $d_G(s,v) > d$.
It has constant query time and a total update time of $O(md+n)$, where $m$ is the number of edges in $G$.
\end{lemma}
}

\paragraph{Even-Shiloach tree (ES-tree).} An ES-tree is a central data structure for maintaining distances in a decremental setting. It maintains the exact distances from a source $s$ under edge deletions, giving the following.

\begin{lemma}[\cite{ES81,henzinger1995fully,king1999fully}]\label{lemma_ES}
There is a data structure called ES-tree that, given a weighted directed graph $G$ undergoing deletions and edge weight increases, a root node $s$, and a depth parameter $d$, maintains, for every node $v$ a value $\ell_s(v)$ such that $\ell_s(v) = d_G(s,v)$ if $d_G(s,v) \leq d$ and $\ell_s(v) = \infty$ if $d_G(s,v) > d$.
It has constant query time and a total update time of $O(md+n)$, where $m$ is the number of edges in $G$.
\end{lemma}

\paragraph{Monotone Even-Shiloach tree (Monotone ES-tree).}
A monotone ES-tree is a generalization of ES-tree studied in \cite{HKN2016,henzinger2014decremental} that handles edge insertions. This is a data structure that maintains distances from a source $s$ in a graph undergoing deletions, insertions and edge weight increases. Note that when edges are inserted to the graph, it can potentially decrease some distances in the graph. However, the data structure ignores such decreases, and maintains distance estimates $\ell_s(v)$ from $s$ to any other node, that can only \emph{increase} during the algorithm. In particular, these distance estimates are not necessarily the correct distances in the graph, and when using this data structure the goal is to prove that these estimates still provide a good approximation for the current distances. We call $\ell_s(v)$ the level of $v$ in the monotone ES-tree rooted at $s$.

We next describe the main properties of monotone ES-trees that we use in our algorithms. A pseudo-code appears in Appendix~\ref{sec:EStrees}. The following lemma follows from \cite{HKN2016} and from the pseudo-code.

\begin{lemma}\label{obs:simple_observations_monotone_ES}
The following holds for a monotone ES-tree of a graph $H$ rooted at $s$, that is maintained up to distance $d$:
\begin{enumerate}[label=(\arabic{*})]
\item At the beginning of the algorithm, $\ell_s(v)=d_H(s,v)$ for all nodes $v$ where $d_H(s,v) \leq d$.\label{obs:beginning}
\item It always holds that $d_H(s,v) \leq \ell_s(v)$.\label{item:easy_side}
\item \label{item: observation one} The level $\ell_s(v)$ of a node $v$ never decreases.
\item When the level $\ell_s(v)$ increases to a value that is at most $d$, we have $\ell_s(v) = \min\limits_{x\in N(v)} \{\ell_s(x)+w_{H}(x,v) \}$.\label{item:change}
\end{enumerate}
\end{lemma}


\begin{lemma}[\cite{HKN2016}]\label{lm:ES-tree}
    For every $d \geq 1$, the total update time of a monotone ES-tree up to maximum level $d$ on a graph $H$ undergoing edge deletions, edge insertions, and edge weight increases is $O(\mathcal E(H)\cdot d+\mathcal W(H)+n)$, where $\mathcal E(H)$ is the total number of edges ever contained in $H$ and $\mathcal W(H)$ is the total number of updates to $H$.
\end{lemma}



\paragraph{Handling large distances.} 

In \cite{HKN14}, the authors show an algorithm that obtains near-additive approximations for decremental APSP, giving the following.

\begin{theorem}[Theorem 1.3 in \cite{HKN14}]\label{thm_large}
For any constant $0 < \epsilon < 1$ and integer $2 \leq k \leq \log{n}$, there is a $(1+\epsilon, 2(1+2/\epsilon)^{k-2})$-approximation algorithm for decremental APSP with constant worst-case query time and expected total update time of $\tilde{O}(n^{2+1/k} (37/\epsilon)^{k-1})$.
\end{theorem}

This will help us in handling \emph{large distances}, because of the following reason. Note that if we have a $(1+\epsilon,\beta)$-approximation algorithm for APSP, then for any pair of nodes at distance $\Omega(\beta/\epsilon)$ from each other, the near-additive approximation is actually a $(1+O(\epsilon))$-approximation because the additive $\beta$ term becomes negligible. Hence, we can focus our attention on estimating distances between pairs of nodes at distance $O(\beta/\epsilon)$ from each other.


\subsection{Warm-up: \texorpdfstring{$(1+\epsilon,2)$}{(1+epsilon,2)}-APSP }\label{sec:additive_warm_up}

Our goal in this section is to use a simpler algorithm based on \cite{DHZ00} that let us maintain $(1+\epsilon,2)$-APSP in total update time of $O(m^{1/2}n^{3/2+o(1)})$ in unweighted undirected graphs with $m=\Omega(n)$ edges. Such an algorithm would outperform the approach in the \Cref{sec:TZ_2+eps} when $m= \Omega( n^{5/3+o(1)} )$, and it also gives a better \emph{near-additive} approximation. We focus on obtaining an algorithm that can maintain $+2$-additive approximate distances for pairs within distance $d$ in total update time $\tilde{O}(m^{1/2}n^{3/2}d)$. This approach can be combined with Theorem~\ref{thm_large} to get our desired bounds. 

\remove{
The static version of the data structure is as follows:

\begin{itemize}
    \item Set $s_1= (\frac{m}{n})^{1/2}$. Let $V_1$ be the set of \textit{dense nodes}: $V_1:= \{ v \in V: \deg(v) \geq s_1\}$. Also let $E_2$ be the set of \textit{sparse edges}, i.e. edges with at least one endpoint with degree less than $s_1$.
    \item Construct a hitting set $D_1$ of nodes in $V_1$. %by adding each node to $D_1$ with probability $\frac{c\ln{n}}{d_1}$ for a constant $c \geq 2$ for our choice. 
    This means that every node $v \in V_1$ has a neighbor in $D_1$. The size of $D_1$ is $O(n \log{n}/s_1)$.
    \item Let $E^*$ be a set of size $O(n)$ such that for each $v \in V_1$, there exists $u \in D_1 \cap N(v)$ such that $\{u,v\} \in E^*$. 
    \item Store distances $D_1 \times V$ by running a BFS from each $u \in D_1$ on the input graph $G$. 
    \item For each $u \in V \setminus D_1$ store a shortest path tree, denoted by $T_u$  rooted at $u$ by running Dijkstra on $(V, E_2 \cup E^* \cup E^{D_1}_u)$, where $E^{D_1}_u$ is the set of weighted edges corresponding to distances in $(\{u\} \times D_1)$ computed in the previous step. 
\end{itemize}
}

\subsubsection{Dynamic Data Structure} 
In \Cref{sec:add_overview}, we described the static algorithm of \cite{DHZ00}. We next describe the dynamic version of the algorithm. As mentioned above, we first focus on maintaining a data structure for pairs of nodes at distance at most $d$ from each other. The algorithm is a variant of the static algorithm. At a high-level, we use ES-trees to maintain distances from nodes in $D_1$ in the graph $G$, and we use monotone ES-trees to maintain distances from nodes $u \in V \setminus D_1$ in the graph $H_u:=(V, E_2 \cup E^* \cup E^{D_1}_u)$. 
The algorithm is summarized as follows.

\begin{enumerate}
    \item \textbf{Node set $D_1$.} We define a random set $D_1$ by sampling each node to $D_1$ with probability $\frac{c\ln{n}}{s_1}$ for $s_1= (\frac{m}{n})^{1/2}$ and a constant $c \geq 2$ of our choice.
    \item \textbf{Edge set $E^*$.} For each node $v$ that has a neighbor in $D_1$, we add one edge $\{v,u\}$ to $E^*$, where $u$ is a neighbor of $v$ in $D_1$.
    \item \textbf{Edge set $E_2$.} For each node $v$ that does not have a neighbor in $D_1$, we add all the adjacent edges to $v$ to the set $E_2$.
    \item \textbf{Maintaining ES-trees from nodes in $D_1$.} For every $u \in D_1$, we maintain an ES-tree rooted at $u$ up to distance $d+2$ in the graph $G$.
    \item \textbf{Maintaining monotone ES-trees from nodes in $V \setminus D_1$.} For each node $u \in V \setminus D_1$, we maintain a monotone ES-tree rooted at $u$ up to distance $d+2$ in the graph $H_u = (V,E_2 \cup E^* \cup E_u^{D_1})$, where $E_u^{D_1}$ are all the edges of the form $\{u,v\}$ where $v \in D_1$. The weight of an edge $\{u,v\} \in E_u^{D_1}$ is $\ell_v(u)$, the level of $u$ in the ES-tree rooted at $v$ computed in the previous step.
\end{enumerate}

%The main challenges in maintaining this data structure decrementally is that 1) there are insertions into the set $E_2$. 2) An edge $(x,y) \in E^*$ may be deleted where before the update $y \in D_1$ and we need to replace it with another edge $(x, y'), y' \in D_1$. This changes $H_u$ which in turn may changes the level and parent of $x$ in $T_u$ and for some $u \in V \setminus D_1$.
%For this we should use a variant of monotone ES-tree from each $u \in V \setminus D_1$ on $E_2 \cup E^* \cup E^{D_1}_u$, and argue that whenever an edge is inserted, we can use use the estimate from the previous update. At a high level, in the monotone ES-tree, whenever the distance of a node $v$ to the root $u$ decreases, we do not change the level (and just update the parent node). We need to argue that distances returned are still within the desired stretch. Intuitively, whenever we use an estimate from a previous update in a tree since distances in $G$ are only increasing, our old estimate is still within the desired stretch after an update. 

\paragraph{Maintaining the monotone ES-trees.}
Next, we give more details about maintaining the monotone trees from nodes $u \in V \setminus D_1$.
%At a high-level, we maintain monotone ES-trees from nodes $u \in V \setminus D_1$, where the input is a dynamic graph $H_u = E_2 \cup E^* \cup E^{D_1}_u$ as follows. Initially these sets of edges are defined as in the static algorithm. As $G$ undergoes edge deletions, the graph $H_u$ will be updated accordingly with respect to the current sets $E_2, E^*, E^{D_1}_u$. Note that an (unweighted) edge deletion in $G$ can result in an edge weight increase in $H_u$.
\begin{enumerate}
    \item Initially, we construct a shortest path tree from $u \in V \setminus D_1$ in the original graph $H_u$, up to distance $d+2$.
    \item Each time the distance estimate between $u$ and a node $x \in D_1$ increases, we update the weight of the edge $\{u,x\} \in E^{D_1}_u$. Note that we can maintain these distances exactly (up to distance $d+2$) as we compute them in the ES-tree rooted at $x \in D_1$.\label{step_dom}
    \item Each time that an edge $\{v,x\} \in E^*$ between a node $v \in V \setminus D_1$ and $x \in D_1$ is deleted from the graph $G$, we distinguish between two cases. If $v$ still has a neighbor $w \in D_1$, we first add a new edge $\{v,w\}$ to the set $E^*$ and then delete the edge $\{v,x\}$.%\sebastian{Why can't we add all such edges to $E^*$ initially? Doesn't the analysis below show that's fine?} 
    Otherwise, we add to $E_2$ all the edges incident to $v$ in the current graph $G$, and then delete the edge $\{v,x\}$.\label{step_E*_ES-tree}
    \item If an edge $e \in E_2$ is deleted, we delete it from $H_u$.
\end{enumerate}
%\ytodo{For running time we should argue that new edges to dominators to each $H_u$ is not causing a problem}

There are several challenges in maintaining the monotone ES-trees. First, during the algorithm there are insertions to the graph $H_u$, as new edges are added to $E_2$ and $E^*$, and our goal is to show that our distance estimates still give a good approximation to the actual distances. In the static version of the algorithm, the algorithm computes the exact distances in the graphs $H_u$ and this provides a good approximation for the distances. However, in our case, the monotone ES-trees are not guaranteed to keep the actual distances in $H_u$ because they ignore distance decreases that come from edge insertions, and we should prove that this does not affect the stretch of the algorithm. 
%Note that there are several challenges in implementing Step~\ref{step_E*} of the algorithm. 

In addition, in the static version of the algorithm it is easy to argue that the graphs $H_u$ are sparse. For example, $E^*$ is clearly a set of at most $O(n)$ edges. In our case, as $H_u$ changes dynamically, we need to prove that the total number of edges added to it is small.
In particular, as each time that an edge in $\{v,x\} \in E^*$ gets deleted, we replace it by another edge, it could be the case that overall in the algorithm we add many different edges to $E^*$ which can affect the total update time of the algorithm. In addition, each time that a node $v$ no longer has a neighbor in $D_1$, we add all its adjacent edges to $E_2$, and we should argue that the number of edges added to $E_2$ is small.
We will discuss these issues in more detail later, and show that we can implement the monotone ES-trees efficiently. 
Next, we provide the stretch analysis of the algorithm.

\subsubsection{Stretch Analysis} \label{sec:add_warmup_stretch} %\mtodo{add figures?} 
For any pair of nodes $u,v$, we denote by $\ell_u(v)$ the level of $v$ in the ES-tree or monotone ES-tree rooted at $u$. We start with the following simple observation. 
%\ytodo{we may wanna make the following 1+eps, as that's what we get if we use emulators/hopsets}
\begin{claim}\label{obs:stretch_warmup}
For any $u \in V, x \in D_1$ such that $d_G(u,x) \leq d+2$, we have that $\ell_x(u)=\ell_u(x)=d_G(u,x).$
\end{claim} 

\begin{proof}
As $x \in D_1$, we maintain the exact distances from $x$ up to distance $d+2$ in the ES-tree rooted at $x$. Hence, we have $\ell_x(u)=d_G(u,x)$. If $u \in D_1$, the same argument holds also for $\ell_u(x)$. If $u \in V \setminus D_1$, in Step~\ref{step_dom} of the algorithm for maintaining the ES-tree we make sure that we also maintain the correct distance between $u$ and $x$ in the monotone ES-tree rooted at $u$, hence $\ell_u(x)=d_G(u,x)$. Note that as the edge $\{u,x\}$ always appears in $H_u$, and by part~\ref{item:change} of Lemma~\ref{obs:simple_observations_monotone_ES} we always maintain the exact distance estimate in this case.
\end{proof}


\begin{claim}\label{claim:stretch_warmup}
For any $u,v \in V$ such that $d_G(u,v) \leq d$, we have that $d_G(u,v) \leq \ell_u(v) \leq d_G(u,v) +2.$
\end{claim}

\begin{proof}
If $v \in D_1$, this already follows from Claim~\ref{obs:stretch_warmup}. Hence we focus on the case that $v \in V \setminus D_1$. The left side of the inequality follows as $d_G(u,v) \leq d_{H_u}(u,v) \leq \ell_u(v)$, using part~\ref{item:easy_side} of Lemma~\ref{obs:simple_observations_monotone_ES}, and since the distances in $H_u$ can only increase compared to distances in $G$. This holds as the edges in $H_u$ are either edges in $G$ or have a weight that is at least the weight of the shortest path between the corresponding endpoints in $G$.
We next show that $\ell_u(v) \leq d_G(u,v) +2$.
We prove the claim by double-induction. First, on the time in the algorithm, and second on the distance $d_G(u,v)$.

\paragraph{Base Case.} At the beginning of the algorithm, this follows from the correctness of the static algorithm, we give the details for completeness. Note that at the beginning of the algorithm, from part~\ref{obs:beginning} of Lemma~\ref{obs:simple_observations_monotone_ES}, we have that $\ell_u(v)=d_{H_u}(u,v)$ as long as $d_{H_u}(u,v) \leq d+2$, and we show that this is indeed the case when $d_G(u,v) \leq d$. 

\begin{itemize}
    \item [Case 1:] $v$ has a neighbor in $D_1$. In this case there is an edge $\{v,x\} \in E^*$ such that $x \in D_1$. Note that as $x$ is a neighbor of $v$ it holds that $d_G(u,x) \leq d+1$. As $x \in D_1$ and $d_G(u,x) \leq d+1$, the graph $H_u$ has an edge $\{u,x\}$ with weight $d_G(u,x)$, and it also has the edge $\{v,x\} \in E^*$. Hence, at the beginning of the algorithm, we have $$\ell_u(v) = d_{H_u}(u,v) \leq d_G(u,x) + d_G(v,x) = d_G(u,x)+1.$$ Now, by the triangle inequality $$d_G(u,x) \leq d_G(u,v)+d_G(v,x)=d_G(u,v)+1.$$ Hence, overall, we get that $\ell_u(v) \leq d_G(u,v) +2$. Note that we get the $d_{H_u}(u,v) \leq d+2$, which implies that the distance is indeed maintained correctly in the monotone ES-tree.
    \item [Case 2:] $v$ does not have a neighbor in $D_1$. In this case, all the edges adjacent to $v$ are in $E_2 \subseteq H_u$. Let $\{w,v\}$ be the edge adjacent to $v$ in the shortest path between $u$ and $v$ in $G$. Then, $d_G(u,w) = d_G(u,v)-1$. From the induction hypothesis, we have that $\ell_u(w) \leq d_G(u,w) +2$. Since the edge $\{w,v\}$ exists in $H_u$, we get that $$\ell_u(v) = d_{H_u}(u,v) \leq d_{H_u}(u,w) + d_{H_u}(w,v)= \ell_u(w)+1 \leq d_G(u,w)+3 = d_G(u,v) + 2.$$ As $d_{H_u}(u,v) \leq d+2$, we indeed maintain the value correctly in the monotone ES-tree.
\end{itemize}


\paragraph{Induction Step.} Assume now that the claim holds after $t-1$ edge deletions, and we will prove it holds after $t$ edge deletions. We prove it by induction on the length of the path between $u$ and $v$ in the current graph $G$. In the base case that $u=v$ it trivially holds.
First, if the value $\ell_u(v)$ did not change after the $t$-th edge deletion, then the claim follows from the induction hypothesis (as the distance $d_G(u,v)$ can never decrease). Hence, we focus on the case that $\ell_u(v)$ changed. By parts~\ref{item: observation one} and~\ref{item:change} of Lemma~\ref{obs:simple_observations_monotone_ES}, when the level $\ell_u(v)$ is changed it is necessarily increased, and the new value is equal to $\min\limits_{x\in N(v)} \{\ell_u(x)+w_{H_u}(x,v) \}$. 

\begin{itemize}
    \item [Case 1:] When the value $\ell_u(v)$ increases, $v$ has a neighbor in $D_1$. By Step~\ref{step_E*_ES-tree} of the algorithm for maintaining the ES-tree, it is guaranteed that there is an edge $\{v,x\} \in E^* \subseteq H_u$. Also, $\ell_u(x)=d_G(u,x)$ by Claim~\ref{obs:stretch_warmup} and since $d_G(u,x) \leq d+1$, as $x$ is a neighbor of $v$. Hence, when we update the value $\ell_u(v)$, the new value is at most $\ell_u(x)+w_{H_u}(x,v)=d_G(u,x)+1.$ The exact same proof that we used in the base case, shows that $d_G(u,x) \leq d_G(u,v)+1$. Hence, $\ell_u(v) \leq d_G(u,v) +2$, as needed. Again, as this value is bounded by $d+2$ it is maintained correctly.
    \item [Case 2:] When the value $\ell_u(v)$ increases, $v$ does not have a neighbor in $D_1$. It follows from Step~\ref{step_E*_ES-tree} of the algorithm for maintaining the ES-tree that all the edges adjacent to $v$ are in $H_u$. Let $x$ be the neighbor of $v$ in the shortest $u-v$ path in $G$. By the induction hypothesis, $\ell_u(x) \leq d_G(u,x)+2$. Hence, $$\ell_u(v) \leq \ell_u(x)+w_{H_u}(x,v) \leq d_G(u,x)+3 = d_G(u,v)+2,$$ this value is again at most $d+2$, hence it is maintained correctly, completing the proof.\qedhere
\end{itemize}
\end{proof}


%We call a node $w$ \textit{stretched} at time $t$ with respect to a Monotone ES-tree $u$ if there is an insertion (potential distance decrease) $(v,w)$ for any $v$ at time $t$.
%\ytodo{old claim, should be removed but we need $(1+\epsilon)$ above}
%\begin{claim}
%For any root $u \in V \setminus D_1$, and any node $v$, we always have $\ell_u(v) %\leq (1+\epsilon)d_G(u,v) +2$.
%\end{claim}
%\begin{proof}[Proof sketch.]
%Consider the shortest path $\pi$ between $u$ and $v$ in $G$. We use an induction on update time $t$. At $t=0$ the claim follows from the static argument. If there are no stretched nodes on $\pi$ after an update at time $t$ then again the claim follows from correctness at time $t-1$ and the static argument. Now assume that there exists a stretched node $u_1$ on path $\pi$. Let $u_1$ be the last such node on $\pi$. Consider two cases:

%\paragraph{Case 1:} Node $u_1$ was stretched because $\deg(u_1)$ becomes less than $d_1$ at time $\leq t$ and hence edges incident to $u_1$ are added to $E_2$. This means at time $t-1$ there was a node $x\in D_1$ adjacent to $u_1$, and thus 
%\[\ell_u(u_1) \leq \ell_u(x) +1 \leq (1+\epsilon)d_G(u,x) + 1 \leq (1+\epsilon)d_G(u,u_1) + 2 \]

%\paragraph{Case 2:} Node $u_1$ is a high-degree node and the edge $(u_1,x) \in E^*$ was removed and replaced with another edge $(u_1, x') \in E^*$ such that $\ell_u(x') \leq \ell_u(x)$. In this case we can use the same argument as before but this time for $x'$. 
%\end{proof}

\subsubsection{Time Complexity} \label{sec:warmup_add_time}

To analyze the time complexity of the algorithm, we first bound the total number of edges added to $E_2$ and $E^*$ in the algorithm. 

\begin{claim} \label{claim_E2}
The total number of edges added to $E_2$ during the algorithm is bounded by $O(ns_1)$ with high probability.%\sebastian{Should we rather go for an in-expectation bound here because we anyway do that for the next claim as well?}
\end{claim}

\begin{proof}
By definition, the edges in $E_2$ are the edges adjacent to all nodes that do not have a neighbor in $D_1$. We show that as long as the degree of a node $v$ is $s_1$ then with high probability it has a neighbor in $D_1$. This follows as $D_1$ is obtained by sampling each node with probability $\frac{c \ln{n}}{s_1}$. Hence, as long as the degree of $v$ is at least $s_1$, the probability that none of its neighbors is in $D_1$ is at most $(1-\frac{c\ln{n}}{s_1})^{s_1} \leq e^{-c\ln{n}}=\frac{1}{n^c}$. It follows that with high probability each node adds less than $s_1$ edges to $E_2$, proving the claim.
\end{proof}

\begin{claim} \label{claim_warmup_E*}
The number of edges added to $E^*$ during the algorithm is bounded by $O( (mn)^{1/2}\log{n} )$ in expectation.
\end{claim}

\begin{proof}
Recall that each node is added to the set $D_1$ with probability $\frac{c \ln{n}}{s_1}$. Hence, in expectation each node $v$ has $O(\deg(v) \frac{\log{n}}{s_1})$ neighbors in the set $D_1$, where $\deg(v)$ is the degree of $v$ at the beginning of the algorithm. These edges are candidates to be added to the set $E^*$, and in the worst case all of them are added to $E^*$ during the algorithm. Summing over all nodes, the expected number of edges added to $E^*$ is bounded by 
\begin{align*}
    \sum_{v \in V} O \left( \frac{\deg(v)\log{n}}{s_1} \right) = O \left( \frac{m\log{n}}{(\frac{m}{n})^{1/2}} \right) = O( (mn)^{1/2}\log{n} ).
\end{align*} 
\end{proof}

\begin{lemma}
The total update time of the algorithm is $\tilde{O}(m^{1/2}n^{3/2}d)$ in expectation.
\end{lemma}

\begin{proof}
Computing the set $D_1$ takes $O(n)$ time by sampling each node to the set independently with probability $\frac{c \ln{n}}{s_1}$. 

Maintaining the sets $E^*$ and $E_2$ is done as follows. At the beginning of the algorithm, each node $v$ scans its neighbors until it finds the first neighbor $u$ in the set $D_1$ if exists, and it adds the edge $\{v,u\}$ to the set $E^*$. If there is no such $u$, then $v$ adds all its adjacent edges to the set $E_2$.
During the algorithm, when the edge $\{v,u\}$ is deleted, $v$ continues scanning its neighbors until it finds the next neighbor $w$ in the set $D_1$ if exists. If $w$ exists it adds the edge $\{v,w\}$ to $E^*$, and otherwise it adds all its adjacent edges to $E_2$. Note that in the whole algorithm, $v$ only needs to scan its neighbors once to find all the edges it adds to $E^*$. Hence, overall maintaining the sets $E^*$ and $E_2$ takes $O(m)$ time. We can also keep a list of the new edges added to $E^*$ and $E_2$ after a deletion, to help maintaining the monotone ES-trees.

We maintain the ES-trees from nodes $u \in D_1$ as follows. Every time that an edge is deleted we delete it from all the trees. In addition, every time that a level $\ell_u(v)$ increases for $u \in D_1, v \in V \setminus D_1$, we update the weight of the edge $\{v,u\}$ in the graph $H_v$, and update it in the monotone ES-tree rooted at $v$.
Maintaining the ES-trees from nodes in $D_1$ takes $O(|D_1|(md+n))$ time by Lemma~\ref{lemma_ES}. By the definition of $D_1$, this equals $$O \left( \frac{n \log{n}}{s_1} \left( md \right) + n^2 \right) =O(m^{1/2}n^{3/2} (\log{n}) d)$$ in expectation, when $m=\Omega(n).$

We maintain the monotone ES-trees rooted at nodes $u \in V \setminus D_1$ as follows. After an edge $e$ is deleted, we update the weights of the edges $\{u,v\} \in E_u^{D_1}$ when their weights are increased in the ES-trees rooted at nodes $v \in D_1$, as described above. In addition, we add to $H_u$ all the edges that were added to $E_2$ and $E^*$ after the deletion of the edge $e$ (we kept them in a list as described above), and finally we delete the edge $e$ from the monotone ES-tree. By Lemma~\ref{lm:ES-tree}, the total update time for maintaining the monotone ES-tree rooted at the node $u$ is 
$O(\mathcal E(H_u)\cdot d+\mathcal W(H_u)+n)$, where $\mathcal E(H_u)$ is the total number of edges ever contained in $H_u$ and $\mathcal W(H_u)$ is the total number of updates to $H_u$. There are $O(n)$ edges in $E_u^{D_1}$ and each of them is updated $O(d)$ times, as the distance estimates only increase and can have $O(d)$ values. This adds $O(nd)$ term to the complexity of one tree, and $O(n^2 d)$ summing up over all trees. The edges in $E_2$ and $E^*$ are only added and removed once in the algorithm. The total number of edges in $E_2$ is $O(ns_1)$ with high probability from Claim~\ref{claim_E2}. Summing up over all trees these edges add $O(n^2 s_1 d)= O(n^{3/2} m^{1/2} d)$ term to the complexity. The total number of edges in $E^*$ is bounded by $O( (mn)^{1/2}\log{n} )$ in expectation from Claim~\ref{claim_warmup_E*}. Summing up over all trees they add $O(n^{3/2} (\log{n}) m^{1/2} d)$ term to the complexity, completing the proof.
\end{proof} 

\begin{claim}
The query time of the algorithm is $O(1)$.
\end{claim}

\begin{proof}
To answer a distance query for a pair of nodes $u,v$, we just return the value $\ell_u(v)$. This can be done in $O(1)$ time if we keep for each node $u$ an array of the values $\ell_u(v)$ which is updated each time a distance estimate change.
\end{proof}

In conclusion, we get the following.

\begin{theorem}
Given an undirected unweighted graph $G$, there is a decremental data structure that maintains $+2$-additive approximation for APSP for pairs of nodes at distance at most $d$ from each other. The expected total update time is bounded by $\tilde{O}(m^{1/2}n^{3/2}d)$, and the query time is constant.
\end{theorem}%\ttodo{Can we not make this whp?}

To get rid of the $d$ term in the complexity, we can combine our approach with Theorem~\ref{thm_large}. This would lead to a $(1+\epsilon,2)$-approximation and replace the $d$ term in the complexity by $n^{o(1)}$ term. We will provide full details of the approach when we describe our general algorithm in the next section (see in particular \Cref{sec_summary_additive}). The high-level idea is to use our algorithm with a parameter $d=n^{o(1)}$, and combine it with Theorem~\ref{thm_large} for computing $(1+\epsilon,\beta)$-approximation for large distances. For an appropriate choice of the parameters, we will get a $(1+\epsilon)$-approximation for the distances between pairs of nodes at distance greater than $d$ from each other, and a $+2$ approximation for pairs of nodes at distance at most $d$ from each other. 

\remove{
%For dealing with the first issue, we note that at any time we only need to keep one edge to a dominator for each node. To apply this change in the monotone ES-trees for all $u \in V \setminus D_1$ in $T_u$ we check whether the parent and level of $u$ needs to be updated due to the replaced edge in $E^*$. Noe that as this is a monotone ES-tree we keep the level of $u$ the same if the update causes a distance decrease. We must now argue that changing all these $n$ trees over all the updates can be done in the desired overall time.
\begin{lemma}\label{lem:update_dominators}
Given updates to the input graph $G=(V,E)$ and a distance bound $d$, the overall time for maintaining $E^*$ and  for all $u \in v\in V \setminus D_1$ $H_u$ and the corresponding (parents and level) changes in $T_u,$ is w.h.p.~$\tilde{O}(m^{1/2}n^{3/2}d)$.
\end{lemma}
\begin{proof}
Each time we exchange an edge in $E^*$ this takes $O(n)$ time to update all $n$ ES-trees $T_{u}$. We bound the total number of candidate edges for $E^*$ to get an upper bound on the total update time needed to make these changes. Suppose $D_1$ is created by sampling each node with probability $\tfrac{\log(n)}{s_1}$, then we get in expectation:
\begin{align*}
    \E\left[ \left|E \cap \left(V \times D_1\right)\right|\right] &= \E\left[ \sum_{v\in V} \left|E \cap \left(\{v\} \times D_1\right)\right|\right] \\
    &= \sum_{v \in V} \E\left[  \left|E \cap \left(\{v\} \times  D_1\right)\right|\right] \\
    &= \sum_{v \in V} \tfrac{\deg(v)\log(n)}{s_1} \\
    &=  \tfrac{m\log(n)}{(\frac{m}{n})^{1/2}} \\
    &= (mn)^{1/2}\log(n).
\end{align*}
So for all $n$ trees, we obtain $O\left( n (mn)^{1/2}\log(n) \right)=\tilde O\left( m^{1/2}n^{3/2} \right) $ for the expected time which can be turned into a high probability bound with a straightforward tail bound.

Now by observing the monotoncity and following the same analysis as the standard ES-tree we get overall time of $\tilde{O}(m^{1/2}n^{3/2}d)$. \ytodo{more details on standard ES-tree argument? }
\end{proof}

Next we argue that we can get the desired time despite the insertions to $E_2$. Roughly speaking this follows since 1- We are maintaining \textit{monotone} ES-trees, and thus we never decrease the distance estimates. 2- For each node $v$, the incident edges $E_v$ need to be added in $E_2$ and thus the trees $T(u), u \in V$ at most once.

\begin{lemma}\label{lem:update_E2}
Given updates to the input graph $G=(V,E)$ and a distance bound $d$, the overall time for maintaining $E_2$ and for all $u \in v\in V \setminus D_1$ $H_u$ and the corresponding (parents and level) changes in $T_u,$ is w.h.p.~$\tilde{O}(m^{1/2}n^{3/2}d)$.
\end{lemma}
\begin{proof}
As discussed edges incident to each node $v$ are added to $E_2$ at most once and only when $\deg(v)\leq s_1$. Overall the updates and all the trees it takes $O(n^2 s_1)=O(n^{3/2}m^{1/2)})$ time to add these edges to the ES-trees. 
Same as before following the standard monotone ES-tree argument for each tree we have total time of $\tilde{O}(n s_1 d)$ and over all the trees we have total time of $O(n^{3/2}m^{1/2)}$.
\end{proof}

\mtodo{combine with time complexity section}
\begin{lemma}\label{lem:warm_up_time}
We can maintain $(1+\epsilon, 2)$-all-pairs distances for pairs of nodes within distance $d$ in total time $O(m^{1/2}n^{3/2}d)$.
\end{lemma}
\begin{proof}
 We can maintain a monotone ES-tree up to depth $d$ in $O(m'd)$ time on a graph with $m'$ edges. In our algorithm, we maintain monotones ES-trees from $\leq n$ sources on subgraphs of size $m'=O((mn)^{1/2})$. This requires $O(m^{1/2}n^{3/2}d)$ total update time.
Moreover, for maintaining distances in $D_1 \times V$ we use an algorithm that maintains $(1+\epsilon)$-MSSP from a set $S$ of sources of polynomial size in $\tilde{O}(|S|m)$ total update time (e.g.~\cite{LN2020}). Here we have $|S|=\tilde{O}(n/m)^{1/2}$, which is $\tilde{O}(m^{1/2}/n^{1/2})$.
\end{proof}\ytodo{we need to use near-linear emulators}
}
\remove{
\subsection{Using sparse emulators for unweighted Graphs}

In unweighted graph one approach for maintaining the above structure  more efficiently is to use near-additive emulators that would let us limit the additive factor $d$ above to $n^{o(1)}$.

\ytodo{First 3 steps can be replaced with just running the above alg for $d=n^{o(1)}$. We should argue somewhere that the $D_1 \times V$ estimate from the emulator are monotone}
\begin{enumerate}
    \item Sample a hitting set $D_1$ for dense nodes with degree $\geq s_1$ and maintain the edge sets $E^*$ and sparse edges $E_2$ as defined above.
    \item Maintain distances $D_1 \times V$ in $\tilde{O}(m^{1/2}/n^{1/2})$ total time using $(1+\epsilon/2)$-MSSP data structures.
    \item Decrementally maintain a $(1+\epsilon/2, n^{o(1)})$-emulator $H$ of size $|H|=n^{1+o(1)}$ \cite{HKN2016} in total time $O(mn^{o(1)})$.\ytodo{Is this the correct citation?}
    \item Maintain $n^{o(1)}$-bounded \textit{monotone ES-trees} $T_u$ up to depth $n^{o(1)}$ for each $u\in V\setminus D_1$ on $E_2\cup E^*\cup \{u\}\times D_1$. Since $|E_2\cup E^*\cup \{u\}\times D_1|= O((mn)^{1/2})$, for all trees this takes times $O(n \cdot n^{o(1)}\cdot (mn)^{1/2})=O(n^{3/2+o(1)}m^{1/2})$. \label{2:monotone_trees}
    \item Maintain APSP on $H$ using Theorem 1.3 of \cite{HKN14}. This takes time $O(n^{2+o(1)})$ and gives us a $(1+\epsilon)$ estimate for all pairs with distance $\geq n^{o(1)}$. \label{2:static_apsp} \ytodo{Does this have a query step? Is it black-box?}
    \item For all pairs $u,v \in V$ set their distance estimate to be the minimum estimate obtained by Steps~\ref{2:monotone_trees} and~\ref{2:static_apsp}.
\end{enumerate}

\paragraph{Analysis Sketch.} For pairs of nodes with distance $n^{o(1)}$ the stretch analysis follows from the same reasoning as in Claim~\ref{claim:stretch_warmup}. For any pair of nodes with distance more than $n^{o(1)}$ the $(1+\epsilon/2, n^{o(1)})$ emulator stretch leads to an overall $1+ \epsilon$ stretch.
The running time analysis follows from~\ref{lem:warm_up_time} by setting $d=n^{o(1)}$.
\sebastian{Is there an argument why we can't have a pure additive approximation (for large distances)?}
\ytodo{We still need to use $(1+\epsilon)$-MSSP, right? Also near-additive emulators}
}

%\sebastian{Is there an argument why we can't have a pure additive approximation (for large distances)?}

%moved above
\iffalse
\paragraph{Dominator edges.} Each time we exchange an edge in $E^*$ this takes $O(n)$ time to update all $n$ ES-trees $T_{u}$. We bound the total number of candidate edges for $E^*$ to get an upper bound on the total update time needed to make these changes. Suppose $D_1$ is created by sampling each node with probability $\tfrac{\log(n)}{s_1}$\ttodo{is this correct?}, then we get in expectation:
\begin{align*}
    \E\left[ \left|E \cap \left(V \times D_1\right)\right|\right] &= \E\left[ \sum_{v\in V} \left|E \cap \left(\{v\} \times D_1\right)\right|\right] \\
    &= \sum_{v \in V} \E\left[  \left|E \cap \left(\{v\} \times  D_1\right)\right|\right] \\
    &= \sum_{v \in V} \tfrac{\deg(v)\log(n)}{s_1} \\
    &=  \tfrac{m\log(n)}{(\frac{m}{n})^{1/2}} \\
    &= (mn)^{1/2}\log(n).
\end{align*}
So for all $n$ trees, we obtain $O\left( n (mn)^{1/2}\log(n) \right)=\tilde O\left( m^{1/2}n^{3/2} \right) $.
\fi




\subsection{\texorpdfstring{$(1+\epsilon, 2k-2)$}{(1+epsilon,2k-2)}-APSP}\label{sec:add_general}

With this warm-up we can move on to the more general algorithm based on \cite{DHZ00} that leads to a larger stretch but a smaller update time. Our goal is to obtain a $(1+\epsilon, 2k-2)$-approximation in $\tilde{O}(n^{2-1/k}m^{1/k})$ time. The algorithm works in  unweighted undirected graphs. As before, we start by describing a data structure that maintains the approximate distances for pairs of nodes at distance at most $d$ from each other, and then in \Cref{sec_summary_additive} we combine it with Theorem \ref{thm_large} to handle the general case.

\remove{
Our goal is to match the following static bound (up to $n^{o(1)}$):
\begin{theorem}[\cite{DHZ00}]
For every $2 \leq k \leq O(\log n)$ there is a $\tilde{O}(n^{2-1/k}m^{1/k})$ time algorithm for computing APSP distance estimates $\hat{d}$ such that for all $u,v \in V$: 
\[d_G(u,v) \leq \hat{d} (u,v) \leq d_G (u,v)+2(k-1). \]

%Moreover, we can show $\hat{d}(u,v) \leq 3d_G(u,v)-2$.
\end{theorem}

The summary of the static algorithm is as follows:
\begin{enumerate}
    \item Set degree thresholds $s_1, s_2, ...,s_{k-1}$, where $d_i=(\frac{m}{n})^{1-i/k} (\log n)^{i/k}$. Let $E_i$ be set of edges incident to a node with degree $\leq d_{i-1}$. 
    \item Sample hitting sets $D_1, D_2, ..., D_k$\ttodo{To do: describe the sampling} and let $E_1^*, \dots, E_k^*$ be sets of edges such that for all $1 \leq i < k$ and each node with degree in $(d_{i-1}, d_i]$ has one edge incident to $D_i$ in $E_i^*$.
    \item For $i=1$ to $k$:
    \begin{itemize}
        \item For every $u \in D_i$ compute SSSP on the graph $G_{u,i}:=(V, E_i \cup E_i^* \cup (\{u\} \times V))$ and update the distance estimates $\hat{d}$.
    \end{itemize}
\end{enumerate}

Note that this algorithm ensures that any node that is not in $E_i$ has degree more than $d_{i-1}$ and thus has neighbor in $D_i$.
}

\subsubsection{Dynamic Data Structure} 
The algorithm is summarized as follows.
%\ytodo{maybe we could give these names: dominators, edges to dominators, light edges, etc. Or better not add new terms?}
\begin{enumerate}
    \item \textbf{Node sets $D_i$.} We define a series of sets $D_1,D_2,...,D_k$ such that for $i=1,...,k-1$, the set $D_i$ is a random set obtained by sampling each node with probability $p_i=\frac{c\ln{n}}{s_i}$, for $s_i=(\frac{m}{n})^{1-i/k} (\ln{n})^{i/k}$, and a constant $c \geq 2$ for our choice. If a node is sampled to several different sets we keep it only in the set with minimum index $i$ among them. Finally, we define $D_k = V \setminus \bigcup_{i=1}^{k-1} D_i.$ Note that by definition, each node is in exactly one set $D_i$.\label{step_Di} 
    \item \textbf{Indices $i_v$ for all $v \in V$.} For a node $v \in V$, we denote by $i_v$ the minimum index $i$ such that $v$ has a neighbor in $D_i$.\label{step_Vi} %Let $V_i = \{v \in V| i_v = i \}.$ 
    \item \textbf{Edge set $E^*$.} For each node $v$, we add one edge $\{v,x\}$ to $E^*$, where $x$ is a neighbor of $v$ in $D_{i_v}$.\label{step_E*}
    \item \textbf{Edge sets $E_i$.} The edges $E_i$ are all the edges $\{u,v\}$ where $\max\{i_u,i_v\} \geq i$. In other words, it is all the edges adjacent to nodes that do not have a neighbor in $D_{i'}$ for $i' < i$.\label{step_Ei}
    \item \textbf{Maintaining monotone ES-trees.}\label{step_ES} The algorithm maintains monotone ES-trees as follows.
    \begin{itemize}
        \item For every $u \in D_1$, we maintain an ES-tree rooted at $u$ up to distance $d+3k$ in the graph $G$. 
    \end{itemize}
    For $i= 2$ to $k$:
    \begin{itemize}
        \item  For every $u \in D_i$, we maintain a monotone ES-tree rooted at $u$ up to distance $d+3k$ in the graph $H_u^i=(V,E_i \cup E^* \cup E_u^i)$, where $E_u^i$ are all the edges $\{u,v\}$, where $v \in D_{i'}$ for $i' < i$. The weight of the edge $\{u,v\}$ is equal to the value $\ell_v(u)$ in the monotone ES-tree rooted at $v$.
    \end{itemize}
\end{enumerate}

Note that we maintain the monotone ES-trees up to distance $d+3k$. As we show in the stretch analysis, this is needed in order to guarantee that we maintain correct estimates for any pair of nodes at distance at most $d$ from each other. Intuitively, as the distance estimate we maintain is at most $d(u,v)+2k$, we would like to maintain correct estimates for pairs of nodes at distance $d+2k$ from each other. Another $k$ term comes from the inductive argument of the stretch analysis. For more details, see \Cref{sec_additive_stretch}.

In \Cref{sec:add_complexity}, we explain in detail how to maintain the data structures efficiently. When we maintain the graphs $H_u^i$ we make sure that the following holds (for a proof, see \Cref{sec:add_complexity}). 

\begin{claim}\label{claim_Hui}
Let $u \in D_i, v \in V$. Throughout the updates, it always holds that $v$ either has an edge $\{v,w\} \in H_u^i$, where $w \in D_{i'}$ for $i' < i$, or all the edges adjacent to $v$ are in $H_u^i$. 
\end{claim}

We are in the first case of the claim when $i_v <i$, and we are in the second case when $i_v \geq i$.

\subsubsection{Stretch Analysis} \label{sec_additive_stretch}

We need the following claims. First, from part~\ref{obs:beginning} of Lemma~\ref{obs:simple_observations_monotone_ES} we have the following.

\begin{claim} \label{claim_beginning}
At the beginning of the algorithm, for any $u \in D_i$, and any $v$ at distance at most $d+3k$ from $u$ in $H_u^i$, we have $\ell_u(v)=d_{H_u^i}(u,v).$
\end{claim}

We next show the following. 

\begin{claim} \label{claim_dist_smaller_dom}
Let $u \in D_i, w \in D_j$ such that $j < i$, then $\ell_u(w) \leq \ell_w(u)$. 
\end{claim}

\begin{proof}
The claim follows as we always have an edge $\{u,w\}$ with weight $\ell_w(u)$ in the graph $H_u^i$, and the value $\ell_w(u)$ is monotonically increasing during the algorithm. Hence, and from part~\ref{item:change} of Lemma~\ref{obs:simple_observations_monotone_ES},  every time that the estimate $\ell_u(w)$ changes it always changes to a value that is at most $\ell_w(u)$.
\end{proof}


%Note that when $\ell_u(v)$ changes, it gets the value $\min_x \{\ell_u(x)+w_{H_u^i}(x,v) \}$ going over the neighbors $x$ of $v$.

We next prove the stretch guarantee of the algorithm.

%\mtodo{try: maintain distances up to $d+3k$?, update algorithm and claims. also relevant for the previous algorithm.}

\begin{lemma}\label{lem_additive_stretch}
For any node $u \in D_i$, and any node $v \in V$ where $d_G(u,v) \leq d+(k-i)$, throughout the updates we have $d_G(u,v) \leq \ell_u(v) \leq  d_G(u,v)+ 2(i-1)$.
\end{lemma}

\begin{proof}
First, the left inequality holds since $d_G(u,v) \leq d_{H_u^i}(u,v) \leq \ell_u(v)$, where the first inequality follows from the definition of $H_u^i$, and the second follows from part~\ref{item:easy_side} of Lemma~\ref{obs:simple_observations_monotone_ES}.
We prove the right inequality by triple-induction, on the value $i$, on the time in the algorithm, and on the distance $d_G(u,v)$.
For $i=1$, the claim follows as we maintain ES-trees up to distance $d+3k$ from the nodes of $D_1$ in the graph $G$. We next assume that it holds for any $i' < i$, and prove that it holds for $i$. The proof is by induction on the time in the algorithm and on the distance $d_G(u,v)$.

\paragraph{Base Case.} We first show that the claim holds at the beginning of the algorithm. From Claim~\ref{claim_beginning}, at the beginning of the algorithm, for any $u,v$ where $d_{H_u^i}(u,v) \leq d+3k$, we have that $\ell_u(v) = d_{H_u^i}(u,v).$ Hence, our goal is to show that $d_{H_u^i}(u,v) \leq d_G(u,v) + 2(i-1)$. Note that if this holds, then since $d_G(u,v) \leq d+(k-i)$, we get that $d_{H_u^i}(u,v) \leq d+3k$, hence the distance is maintained in the monotone ES-tree as needed.
The proof is by induction on $d_G(u,v)$.

\begin{itemize}
    \item [Case 1:] $i_v < i$. In this case, there is an edge $\{v,w\} \in E^* \subseteq H_u^i$ such that $w \in D_{i'}$ for $i' < i$. Note that since $d_G(u,v) \leq d+(k-i)$, and $w$ is a neighbor of $v$, then we have $d_G(u,w) \leq d +(k-i)+1 = d+(k-(i-1)) \leq d+(k-i')$. Hence, we can use the induction hypothesis on $w$ and $u$. From the induction hypothesis and Claim~\ref{claim_beginning}, we have 
    $\ell_w(u) = d_{H_u^i}(w,u) \leq d_G(w,u)+2(i-2)$ (note that this value is smaller than $d+3k$, so we can indeed use Claim~\ref{claim_beginning}). 
    Hence, we get that $$d_{H_u^i}(u,v) \leq d_{H_u^i}(u,w) + d_{H_u^i}(w,v) \leq d_G(w,u) + 2(i-2)+1.$$ We have, $d_G(w,u) \leq d_G(w,v)+d_G(v,u)=d_G(u,v)+1$.
    Hence, we get, $d_{H_u^i}(u,v) \leq d_G(u,v)+2(i-1)$, as needed.
    %Moreover, from Claim~\ref{claim_dist_smaller_dom}, $\ell_u(w) \leq \ell_w(u)$. At the beginning of the algorithm, we compute a shortest path tree from $u$ in the graph $H_u^i$, hence $\ell_u(v)=d_{H_u^i}(u,v)$.\mtodo{should have some claim} This gives $$\ell_u(v) = d_{H_u^i}(u,v) \leq d_{H_u^i}(u,w) + d_{H_u^i}(w,v) = \ell_u(w) + 1 \leq d_G(w,u) + 2(i-2)+1.$$ We have, $d_G(w,u) \leq d_G(w,v)+d_G(v,u)=d_G(u,v)+1$. Hence, we get, $\ell_u(v) \leq d_G(u,v)+2(i-1)$.
    \item [Case 2:] $i_v \geq i$. In this case, by definition, all the edges adjacent to $v$ are in $E_i \subseteq H_u^i$. Let $w$ be the neighbor of $v$ on the shortest $u-v$ path. From the induction hypothesis, $d_{H_u^i}(u,w) \leq d_G(u,w)+2(i-1)$. So we have $$d_{H_u^i}(u,v) \leq d_{H_u^i}(u,w)+d_{H_u^i}(w,v) \leq d_G(u,w)+1+2(i-1)=d_G(u,v)+2(i-1).$$
\end{itemize}


\paragraph{Induction Step.} We assume that the claim holds after $t-1$ edge deletions, and show that it holds after $t$ edge deletions. The proof is by induction on the distance $d_G(u,v)$. If $u=v$, the claim trivially holds. If the estimate $\ell_u(v)$ did not change after the $t$'th edge deletion, then the claim holds from the induction hypothesis. Note that an edge deletion can potentially lead to several changes of the value $\ell_u(v)$. It is enough to prove that the stretch guarantee holds in the last time the value $\ell_u(v)$ increases, as the value stays the same afterwards. We divide into cases based on Claim~\ref{claim_Hui}.

\begin{itemize}
    \item [Case 1:] When the value $\ell_u(v)$ increases in the last time, there is an edge $\{v,w\} \in H_u^i$ such that $w \in D_{i'}$ for $i' < i$. 
    %$v$ has a neighbor $w$ in $D_{i'}$ for $i' < i$ in the graph $H_u^i$.\\
    From part~\ref{item:change} of Lemma~\ref{obs:simple_observations_monotone_ES}, when $\ell_u(v)$ changes, it gets the value $\min\limits_{x\in N(v)} \{\ell_u(x)+w_{H_u^i}(x,v) \}$, assuming this value is at most $d+3k$, and we will show that this is indeed the case. Hence, $$\ell_u(v) \leq \ell_u(w)+w_{H^i_u}(w,v) \leq \ell_w(u) + 1 \leq d_G(u,w)+2(i-2)+1.$$ Here we use Claim~\ref{claim_dist_smaller_dom} and the fact that $w \in D_{i'}$ for $i' < i$ (note that since $d_G(u,v) \leq d+(k-i)$, and $w$ is a neighbor of $v$, it follows that $d_G(u,w) \leq d+(k-i')$, so we can indeed use the induction hypothesis). It also holds that $d_G(u,w) \leq d_G(u,v)+d_G(v,w)=d_G(u,v)+1$. Hence, overall, we have $\ell_u(v) \leq d_G(u,v) + 2(i-2)+2=d_G(u,v)+2(i-1)$, as needed. Note that this value is indeed smaller than $d+3k$, hence we maintain the correct estimate in the monotone ES-tree.
    \item [Case 2:] When the value $\ell_u(v)$ increases in the last time, all the adjacent edges to $v$ are in the graph $H_u^i$.
    Let $w$ be the neighbor of $v$ on the shortest $u-v$ path. When the distance estimate of $v$ changes we have $$\ell_u(v) \leq \ell_u(w)+d_{H_u^i}(w,v) \leq d_G(u,w) + 2(i-1) + 1 = d_G(u,v)+2(i-1),$$ where we use the induction hypothesis on $w$ and part~\ref{item:change} of Lemma~\ref{obs:simple_observations_monotone_ES}. Note that it follows that the values $\ell_u(v),\ell_u(w)$ are smaller than $d+3k$, hence they are maintained correctly in the monotone ES-tree.\qedhere
\end{itemize}
\end{proof}

%\mtodo{discuss the combination with near-additive APSP}

\iffalse
\begin{proof}[Proof sketch.]\ytodo{needs induction on distances rather than $i$. Do we need double or triple induction?}
Note that the weight of edges $u \times D_i, 1 \leq i < k$ are obtained by a monotone ES-tree on $G_{u,i}$.
We use a double induction on $t$ and distance $d$ from $u$ and induction on $i$. At $t=0$ claim follows from the static analysis in \cite{DHZ00}. Now assume that at time $t-1$ we have correctly maintained the estimates. At time $t$ for $d=1$, the argument is the same as Claim~\ref{claim:stretch_warmup}. Now assume that the claim holds for all $y \in D_{i-1}, x \in V$. Consider any $u \in D_i, v \in V$, and let $\pi$ be the shortest path  between $u $ and $v$ in $G$. If all the edges on $\pi$ are in $E_i$ and there are no insertions/distance decreases along this path then the claim follows from the induction hypothesis on time $t-1$.

The difficult case is when not all edges belong to $E_i$. Then there must be a node $u'$ on $\pi$ and an edge $(u',w_{i-1}) \in E^*$ where $w_{i-1} \in D_{i-1}$. Assume without loss of generality that $w_{i-1}$ is the last such node (with a neighbor in $D_i$). Same as before we can argue that if node $w_i$ is stretched in $T_u$ the estimate comes from time $t-1$ in the tree rooted at $w_i$. We have,
\[ \ell_u(u') \leq \ell_u(w_i)+1 \leq d_G(u,w_i)+1 \leq d_G(u,u')+2 \leq 2(d-2)+2= 2(d-1) \]
\end{proof}
\fi
%Hence overall we get an additive term $2(k-1)$. Same as before such an algorithm would have a dependence on a maximum distance bound $d$ between pairs, and we would use a $(1+\epsilon, \beta)$-emulator to bound this depth to $\beta$. This would add a multiplicative factor to estimates computed. Overall by using an emulator with $\beta=n^{o(1)}$ we get,
%\ytodo{Can we get decremental emulators with more general tradeoffs ($\beta=$ polylog?) If so there is a better balancing}
%\begin{theorem}
%Given an undirected unweighted graph there is a decremental data structure that maintains $(1+\epsilon, 2(k-1))$-APSP with total update time $O(n^{2-1/k+o(1)}m^{1/k})$.
%\end{theorem}

%\ytodoi{I am unsure if the $3$-approx analysis also works in decremental settings.}

%\ytodo{Can we add some additional structure for the $3$-approx, such as adding an edge from each node to all dominators (having $E^*_i$)}

\iffalse
\subsection{Weighted APSP with Hopsets}
For unweighted graphs we can only maintain the emulators if \cite{HKN2016}. If we turn to a weighted version of the algorithm e.g. by Cohen and Zwick \cite{TBA}, then we will need to use decremental hopsets.

We first maintain a $(\beta, \epsilon)$ hopset of size $O(n^{3/2})$ with $\text{polylog}(n)$ hopbound \cite{LN2020}, and use this to maintain each of the above monotone ES-trees in $\tilde{O}(m'+n^{3/2})$ time. In maintaining these trees we need the following properties of the hopset data structure $H$:
\ytodo{Is this the hopset tradeoff we want or one with $\beta=n^{o(1)}$?}
\begin{itemize}
    \item Between every pair of nodes in $G$ there is a path with polylogarithmic hopbound and stretch $(1+\epsilon)$ in $H$.
    \item The hopset edge weights never decrease, in other words distance in $G \cup H$ are non-decreasing.
\end{itemize}

Additionally, in order to use such a hopset, we need to combine it with a standard scaling algorithms that let us maintain $\beta$ hop-sensitive single-sources distances on each monotone ES-tree $T_u$. 
\fi

\subsubsection{Time Complexity} \label{sec:add_complexity}

We next analyze the time complexity of the algorithm. We start by showing that the first steps of the algorithm can be maintained efficiently. 

\begin{claim}
The first 4 steps of the algorithms can be maintained in $O((m+n)k)$ time.
\end{claim}

\begin{proof}
\textbf{Step~\ref{step_Di}.} Computing the sets $D_i$ takes $O(kn)$ time, by just sampling each node to the set $D_i$ with probability $p_i$. For each node $v$, we can stop the sampling after the first time $v$ is sampled to one of the sets. This is only done once at the beginning of the algorithm. For each node $v$, we keep the index $i$ such that $v \in D_i$. For each $i$, we create a list of all the nodes in the set $D_i$. Creating all these lists takes $O(n)$ time by scanning the nodes.

\textbf{Steps~\ref{step_Vi} and~\ref{step_E*}.} The indices $i_v$ and the set $E^*$ are maintained dynamically in the algorithm, as follows. At the beginning, each node $v$ scans its neighbors to find the minimum index $i$ such that $v$ has a neighbor in $D_i$, it also saves the name of an edge $\{v,u\}$ to the first neighbor of $v$ in $D_{i_v}$, this edge is added to set $E^*$. In the algorithm, if the edge $\{v,u\}$ is deleted, then $v$ goes over its next neighbors, and finds the next neighbor $w \in D_{i_v}$ if exists. If such a neighbor exists, then $\{v,w\}$ is added to $E^*$. Otherwise, the value $i_v$ increases, and $v$ scans its neighbors again to determine the new index $i_v$ and adds the first edge $\{v,w\}$ from $v$ to a neighbor in $D_{i_v}$ to $E^*$. Note that overall in the algorithm each node needs to scan its neighbors $O(k)$ times, because as long as the value $i_v$ stays the same, we just need to go over the neighbors once during the whole algorithm to determine all the edges added to $E^*$. Overall these steps take $O(mk)$ time. 

\textbf{Step~\ref{step_Ei}.} Maintaining the sets $E_i$ is done as follows. Initially, we add to $E_i$ all the edges adjacent to nodes $v$ where $i_v \geq i$. During the algorithm, every time that the value $i_v$ of a node $v$ changes, say from $j$ to $j'$, then for all $j < i \leq j'$, we add the edges adjacent to $v$ to the set $E_i$. We also keep a list of all the newly added edges in $E_i$ after a deletion of an edge, so that we can update them in all the relevant monotone ES-trees. The time complexity for maintaining the set $E_i$ is proportional to the number of edges added to $E_i$ during the algorithm, which can trivially be bounded by $O(m)$, hence this step takes at most $O(mk)$ time.  
\end{proof}

\paragraph{Maintaining the monotone ES-trees.} 

The main challenge in the algorithm is maintaining the monotone ES-trees in Step~\ref{step_ES} of the algorithm. Before analyzing the time complexity of the algorithm, we start by describing the algorithm in more detail. Note that each time an edge is deleted it can lead to many changes, including edge insertions and weight changes in all the trees we maintain. We next explain in detail in which order we make these changes in order to maintain the correctness of the algorithm. We use the algorithms for ES trees and monotone ES trees described in \Cref{sec:additive_tools}.

%The main part of the algorithm is maintaining the monotone ES-trees. This is done as follows. 
\subparagraph{ES-trees rooted at nodes of $D_1$.} 
First, for any node $u \in D_1$, we start by creating an ES-tree rooted at $u$ in the graph $G$, maintained up to distance $d+3k$. After creating this tree, for any $v \in D_i$ for $i > 1$, we add an edge $\{u,v\}$ with weight $\ell_u(v)$ to the set of edges $E_v^i$. During the algorithm, every time an edge is deleted, we delete it from the ES-tree and update it accordingly. For any $v$ where $\ell_u(v)$ has increased after the deletion and where $v \in D_i$ for $i > 1$, we update the corresponding edge in $E_v^i$, and save it in a list of edges in $E_v^i$ that were updated after the deletion. This completes the description for $u \in D_1$. 
%\ytodo{We have a monotone ES tree pseudocode right now that in the end that is not consist with the algorithms here. We should either have a standard ES tree there and explain how to make it monotone here or change that algorithm and explain the changes here}
\subparagraph{Monotone ES-trees rooted at nodes of $D_i$.}
For $u \in D_i$ for larger values of $i$, we maintain the monotone ES-trees according to the order of the index $i$, from 2 to $k$. We work as follows. Initially, we create a monotone ES-tree rooted at $u$ up to distance $d+3k$ in the graph $H_u^i=(V,E_i \cup E^* \cup E_u^i)$, with the edges $E_i,E^*$ at the beginning of the algorithm, and with the edges $E_u^i$ that were obtained when creating the monotone ES-trees in previous levels. Note that all these edges of $E_u^i$ are of the form $\{u,v\}$ where $v \in D_{i'}$ for $i' < i$, so we indeed already created the monotone ES-tree rooted at the relevant nodes $v$.  After creating the tree, for any node $v \in D_{i'}$ for $i' > i$, we add a new edge to the set $E_v^{i'}$ with weight $\ell_u(v)$.
During the algorithm, when an edge $e$ is deleted, we update the monotone ES-trees in several steps as follows, again we work according to the order of the index $i$.
%\ytodo{we should explain why this is not black-box of prev monotone ES-tree algorithms (since we are updating $E^*$). In general we may want to make it more clear what is new and what is known from ES trees}
\begin{itemize}
    \item \textit{First step: adding the new edges of $E^*$ and $E_i$ to $H_u^i$.} The deletion of the edge $e$ can lead to adding new edges to $E^*$ or $E_i$ in steps~\ref{step_E*} and~\ref{step_Ei} of the algorithm. As explained in these steps, we maintain the newly added edges, and can add them to a list of newly added edges (to $E^*$ and to $E_i$ separately). When we update the monotone ES-tree rooted at $u \in D_i$, we start by adding all these newly added edges that were added to $E^*$ and $E_i$ to $H_u^i$.
    \item \textit{Second step: updating the edges in $E_u^i$.} Next, there are possibly edges in $\{u,v\} \in E_u^i$ that their weight is increased after the deletion. This is already observed when we update the monotone ES-tree rooted at $v$, and if it happens, we add the edge $\{u,v\}$ with new weight $\ell_v(u)$ to a list of edges that were updated in $E_u^i$. When we update the monotone ES-tree rooted at $u$, we update the weights of all these edges (we already have their list after we finish maintaining the ES-trees in previous levels).
    \item \textit{Third step: deleting the edge $e$.} Lastly, we delete the edge $e$ that was deleted in this step from the monotone ES-tree. 
\end{itemize}

Note that the weight increases and deletion can increase the levels $\ell_u(v)$ of some nodes in the tree. For any $v$ in $D_j$ for $j > i$ such that $\ell_u(v)$ increases, when the value increases, we add an edge $\{v,u\}$ with the new weight $\ell_u(v)$ to the list of newly updated edges in $E_v^j$. This completes the description of the algorithm.
From the above description, we can now prove Claim~\ref{claim_Hui}. 

\begin{proof}[Proof of Claim~\ref{claim_Hui}]
Let $u \in D_i, v \in V$. At the beginning of the algorithm, if $i_v < i$, then there exists an edge $\{v,w\} \in E^* \subseteq H_u^i$ where $w \in D_{i'}$ for $i' < i$. Otherwise, $i_v \geq i$, and all the edges adjacent to $v$ are in $E_i \subseteq H_u^i$. During the algorithm, every time an edge $e$ of the form $\{v,w\} \in E^*$ is deleted, then if $i_v < i$ after the deletion, we add a new edge $\{v,w'\}$ to $E^*$ for $w' \in D_{i'}$ for $i' < i$, and if $i_v \geq i$, we add all the edges adjacent to $v$ to $E_i \subseteq H_u^i$. In the algorithm, we first add these new edges to $H_u^i$, and only then delete the edge $e$, which proves the claim.
\end{proof}

To bound the time complexity, we need the following useful claims that bound the total number of edges added to the sets $E_i$ and $E^*$ during the algorithm.

\begin{claim}\label{claim_Ei}
For $i \geq 2$, the total number of edges added to $E_i$ during the algorithm is bounded by $O(ns_{i-1})$ with high probability.
\end{claim}

\begin{proof}
The edges $E_i$ are all the edges adjacent to nodes $v$ where $i_v \geq i$. We show that at the first time in the algorithm that $i_v \geq i$, the degree of $v$ is at most $s_{i-1}$ with high probability. This implies that each node $v$ adds at most $s_{i-1}$ edges to the set $E_i$, proving the claim. 
To bound the degree of $v$, note that by definition the set $D_{i-1}$ is obtained by sampling each node with probability $\frac{c \ln{n}}{s_{i-1}}$, all these nodes are added to the set $D_{i-1}$, unless they are already part of a set $D_{i'}$ for $i' < i-1$. As long as the degree of $v$ is at least $s_{i-1}$, the probability that none of the neighbors of $v$ is sampled to $D_{i-1}$ is at most $(1-\frac{c\ln{n}}{s_{i-1}})^{s_{i-1}} \leq e^{-c\ln{n}}=\frac{1}{n^c}$. Hence, with high probability $v$ has a neighbor sampled to $D_{i-1}$, implying that $i_v \leq i-1$ as long as $\deg(v) \geq s_{i-1}$. From union bound, this holds for all nodes with high probability. This proves that the total number of edges added to $E_i$ in the algorithm is bounded by $O(ns_{i-1})$. Note that the proof assumes that we have an oblivious adversary and that the order in which edges are deleted does not depend on the random choices of the algorithm.
\end{proof}

\begin{claim}\label{claim_E*}
In expectation, the total number of edges added to the set $E^*$ in the algorithm is bounded by $O(kn^{1-1/k}m^{1/k}(\log{n})^{1-1/k})$.
\end{claim}

\begin{proof}
We analyze the number of edges $v$ adds to the set $E^*$ when $i_v=i$ for all $i$. First, when $i_v=1$, then in the worst case $v$ adds to $E^*$ edges to all its neighbors in $D_1$. By the definition of $D_1$, each node is added to $D_1$ with probability $\frac{c\ln{n}}{s_1}$. Hence, in expectation $v$ has $O(\frac{\deg(v) \log{n}}{s_1})$ neighbors in $D_1$, where $\deg(v)$ is the degree of $v$ at the beginning of the algorithm. When we sum up over all $v$, this adds at most $O(\frac{m \log{n}}{s_1})=O(n^{1-1/k}m^{1/k}(\log{n})^{1-1/k})$ edges in expectation, where the last equality follows from the definition of $s_1$.

When $i_v=i$ for $i > 1$, as we proved in Claim~\ref{claim_Ei}, with high probability the degree of $v$ is at most $s_{i-1}$. Note that each one of the neighbors of $v$ is added to $D_i$ with probability $\frac{c\ln{n}}{s_i}$. Hence, in expectation, $v$ has $O(s_{i-1} \cdot \frac{\ln{n}}{s_i})$ neighbors in $D_i$. In the worst case, it adds edges to all of them during the algorithm. In total, the expected number of edges added by all nodes to $E^*$ in the case that $i_v \geq 2$ is bounded by $\sum_{i=2}^{k} O(\frac{\ n \log{n} s_{i-1}}{s_i})$.
By the definition of $s_i$, we have that $\frac{s_{i-1}}{s_i}=(\frac{m}{n})^{1/k}(\ln{n})^{-1/k}$. Hence, in total the expected number of edges added is bounded by $O(k n^{1-1/k}m^{1/k}(\log{n})^{1-1/k}).$
\end{proof}

Next, we prove the complexity of the algorithm. 

\begin{lemma}\label{lemma_time_additive}
Maintaining the ES-trees takes $O(kn^{2-1/k}m^{1/k}(\log{n})^{1-1/k} d)$ time in expectation.
\end{lemma}

\begin{proof}
Maintaining the ES-trees rooted at nodes in $D_1$ takes $O(|D_1|(md+n))$ time by Lemma~\ref{lemma_ES}. 
Since $|D_1| = O(\frac{n\log{n}}{s_1})$ in expectation, and by the choice of $s_1$, this adds $O(n^{2-1/k}m^{1/k}(\log{n})^{1-1/k} d)$ term to the complexity, when $m=\Omega(n)$. 

From Lemma~\ref{lm:ES-tree}, maintaining a monotone ES-tree rooted at $u \in D_i$ takes $O(\mathcal E(H_u^i)\cdot d+\mathcal W(H_u^i)+n)$ time, where $\mathcal E(H_u^i)$ is the total number of edges ever contained in $H_u^i$ and $\mathcal W(H_u^i)$ is the total number of updates to $H_u^i$. To complete the proof, we should analyze these values. Recall that $H_u^i=(V, E_i \cup E^* \cup E_u^i)$. Note that $E_u^i$ has at most $n-1$ edges, that are updated $O(d)$ times during the algorithm, as we maintain distances up to distance $d+3k$, hence these edges add $O(nd)$ term to the complexity of maintaining one tree, and $O(n^2 d)$ for all $n$ trees. The edges in $E^*$ and $E_i$ are only added and removed once from a tree during the algorithm. From Claim~\ref{claim_Ei}, the total number of edges added to $E_i$ during the algorithm is bounded by $O(ns_{i-1})$ with high probability. Since these edges are part of all the ES-trees rooted at nodes in $D_i$, maintaining these edges adds $O(|D_i| n s_{i-1} d)$ term to the complexity.
This equals $$O \left( \frac{n \log{n}}{s_i} \cdot n s_{i-1} \cdot d \right) = O( n^{2-1/k}m^{1/k}(\log{n})^{1-1/k} d),$$ following the same calculations done in Claim~\ref{claim_E*} (but with an extra $nd$ term). Summing over all $k$ adds another multiplicative $k$ term to the complexity.
From Claim~\ref{claim_E*}, the total number of edges added to $E^*$ is bounded by $O(kn^{1-1/k}m^{1/k}(\log{n})^{1-1/k})$ in expectation. Summing over all $n$ trees, it adds another $O( kn^{2-1/k}m^{1/k}(\log{n})^{1-1/k} d)$ term to the complexity, completing the proof.
\end{proof}

\begin{claim}
The query complexity of the algorithm is constant.
\end{claim}

\begin{proof}
To answer a distance query for a pair of nodes $u,v$, we work as follows. If $u \in D_i, v \in D_j$, and $i<j$, our distance estimate is $\ell_u(v)$, otherwise our distance estimate is $\ell_v(u)$. Note that it is easy to maintain the values $\ell_u(v)$ in a way that the query time is constant. For example, we can maintain for each node $u$ an array with the values $\ell_u(v)$, and update the value in constant time each time it was changed in the algorithm.
\end{proof}


%$O(\mathcal E(H)\cdot d+\mathcal W(H)+n)$, where $\mathcal E(H)$ is the total number of edges ever contained in $H$ and $\mathcal W(H)$ is the total number of updates to $H$.

\subsubsection{Putting Everything Together} \label{sec_summary_additive}

In the last sections we showed an algorithm that for any pair of nodes at distance at most $d$ maintains a $+2(k-1)$-additive approximation in $\tilde{O}(n^{2-1/k}m^{1/k}d)$ time. To get our final result, we combine it with Theorem~\ref{thm_large} to handle large distances, proving the following.

\thmadd*


Note that for small values of $k$ and if $m=n^{1+\rho}$ for a constant $\rho>0$, we get a complexity of $\tilde{O}(n^{2-1/k}m^{1/k})$, where in the general case we have an extra $n^{o(1)}$ term.

\begin{proof}
We maintain simultaneously our dynamic data structure to handle short distances of at most $d$, and the data structure from Theorem~\ref{thm_large} to handle large distances of more than $d$, for a parameter~$d$ which we specify next. We run the algorithm from Theorem~\ref{thm_large} with parameters $\epsilon' = \epsilon/2$, and $k'$ to be specified later. Recall that this algorithm gives a $(1+\epsilon', \beta)$-approximation for $\beta= 2(1+2/\epsilon')^{k'-2}$. We define $d=\beta/\epsilon'$.

To answer a distance query for a pair of nodes $u,v$, we return the minimum between the values obtained in both algorithms. As in both algorithms the query time is constant, we get a constant query time. 

For the stretch analysis, if $d(u,v) \leq d$, then by Lemma~\ref{lem_additive_stretch}, we get an additive stretch of at most $+2(k-1)$. If $d(u,v) \geq d = \beta/\epsilon'$, then the distance estimate we get from the $(1+\epsilon',\beta)$-approximation algorithm is at most $$(1+\epsilon')d(u,v) + \beta \leq (1+\epsilon')d(u,v) + \epsilon' d(u,v) = (1+2\epsilon')d(u,v)=(1+\epsilon)d(u,v).$$
Hence, overall we get a $(1+\epsilon,2(k-1))$-approximation, as needed.

We now analyze the time complexity. From Lemma~\ref{lemma_time_additive}, the time complexity of our algorithm is $O(kn^{2-1/k}m^{1/k}(\log{n})^{1-1/k} d)$, and from Theorem~\ref{thm_large} the time complexity of the near-additive approximation algorithm is $\tilde{O}(n^{2+1/k'} (37/\epsilon)^{k'-1})$.
We distinguish between two choices for the value $k'$. First we choose $k' = k/\rho$, where $m=n^{1+\rho}$. In this case, the complexity of both expressions is bounded by $\tilde{O}((n^{2-1/k}m^{1/k}) O(1/\epsilon)^{k/\rho})$, as $d=\beta/\epsilon', \beta= 2(1+2/\epsilon')^{k'-2}$, and $O(n^{2-1/k}m^{1/k})=O(n^{2-1/k}n^{(1+\rho)/k})=O(n^{2+\rho/k})=O(n^{2+1/k'})$.

Second, we choose $k' = \sqrt{\log{n}}$. In this case, $\tilde{O}(n^{2+1/k'} (37/\epsilon)^{k'-1})= O(n^{2+o(1)})$. Similarly, $d= \beta/\epsilon' = O(1/\epsilon')^{k-1}= O(1/\epsilon)^{O(\sqrt{\log{n}})}=n^{o(1)}$. Hence, in total we get a complexity of $O(n^{2-1/k+o(1)}m^{1/k})$, when $m=\Omega(n)$.
\end{proof}

%For any constant $0 < \epsilon < 1$ and integer $2 \leq k \leq \log{n}$, there is a $(1+\epsilon, 2(1+2/\epsilon)^{k-2})$-approximation algorithm for decremental APSP with constant worst-case query time and expected total update time of $\tilde{O}(n^{2+1/k} (37/\epsilon)^{k-1})$.



\remove{
\paragraph{Dominator edges.} Each time we exchange an edge in $E_i^*$ this takes $O(n)$ time to update all $n$ ES-trees on $G_{u,i}$. We bound the total number of candidate edges for $E_i^*$ to get an upper bound on the total update time needed to make these changes. Suppose $D_i$ is created by sampling each node with probability $\tfrac{\log(n)}{d_i}$\ttodo{is this correct?}, then we get in expectation:
\begin{align*}
    \E\left[ \left|E \cap \left(V \times D_i\right)\right|\right] &= \E\left[ \sum_{v\in V} \left|E_{i+1} \cap \left(\{v\} \times D_i\right)\right|\right] \\
    &= \sum_{(u,v)\in E_i} \E\left[  u\in D_i\right] \\
    &= \sum_{(u,v)\in E_i} \tfrac{\log(n)}{d_i} \\
    &= |E_i|\tfrac{\log(n)}{d_i} \\
    &\leq nd_i\tfrac{\log(n)}{d_i} \\
    &= n\log(n)
\end{align*}
So for all $n$ trees, we obtain $O\left(n^2\log(n)\right)=\tilde O\left(n^2\right) $  time. And in total for all $i=1,\dots k-1$ we obtain total update time $\tilde O\left(kn^2\right) $, which is dominated by $\tilde{O}(n^{2-1/k}m^{1/k})$.

\paragraph{Decremental implementation.} Similar to the warm-up algorithm we use Monotone ES-trees for maintaining distances in the final loop. We note that here the distance estimates $\hat{d}$ for larger iterations are updated based on the smaller estimates, which introduces \textit{insertions} in our trees. Same insight as before but a somewhat more involved induction is needed to prove correctness. We first give a sketch of the correctness analysis and then use similar ideas as before to get an improved update time by using near-additive emulators.
}